\documentclass[10pt,a4paper]{amsart}
\usepackage{amsmath,amssymb,amsthm,graphicx}
% Retained preamble for the Stacks Project chapter {CH} (shared project preamble).
% The chapter file pinned-original-source.tex only carries \input{{preamble}} on its first line,
% so no preamble text lives inside the chapter itself.  The Stacks Project uses one single shared
% preamble for every chapter: the snapshot stages it three times (sources/stacks-chow/preamble.tex,
% sources/stacks-etale-cohomology/preamble.tex and focused-seeds/AGP/sources/stacks-snapshot-20261011/
% preamble.tex) and those copies are identical in size.  The block below is copied mechanically,
% line for line, from sources/stacks-chow/preamble.tex lines 14-18 and 154-217: the Xy-pic setup and
% the native theorem declarations together with the mathematical macros.  Nothing was added, removed
% or redefined; the wrapper class and the amsmath/amssymb/amsthm/graphicx base packages are supplied
% by the assembler.
% For commutative diagrams we use Xy-pic
% Chapter: stacks-spaces-more-morphisms.tex -- More on Morphisms of Spaces.
\usepackage[all]{xy}

% We use 2cell for 2-commutative diagrams.
\xyoption{2cell}

% Theorem environments.
%
\theoremstyle{plain}
\newtheorem{theorem}[subsection]{Theorem}
\newtheorem{proposition}[subsection]{Proposition}
\newtheorem{lemma}[subsection]{Lemma}

\theoremstyle{definition}
\newtheorem{definition}[subsection]{Definition}
\newtheorem{example}[subsection]{Example}
\newtheorem{exercise}[subsection]{Exercise}
\newtheorem{situation}[subsection]{Situation}

\theoremstyle{remark}
\newtheorem{remark}[subsection]{Remark}
\newtheorem{remarks}[subsection]{Remarks}

\numberwithin{equation}{subsection}

% Macros
%
\def\lim{\mathop{\mathrm{lim}}\nolimits}
\def\colim{\mathop{\mathrm{colim}}\nolimits}
\def\Spec{\mathop{\mathrm{Spec}}}
\def\Hom{\mathop{\mathrm{Hom}}\nolimits}
\def\Ext{\mathop{\mathrm{Ext}}\nolimits}
\def\SheafHom{\mathop{\mathcal{H}\!\mathit{om}}\nolimits}
\def\SheafExt{\mathop{\mathcal{E}\!\mathit{xt}}\nolimits}
\def\Sch{\mathit{Sch}}
\def\Mor{\mathop{\mathrm{Mor}}\nolimits}
\def\Ob{\mathop{\mathrm{Ob}}\nolimits}
\def\Sh{\mathop{\mathit{Sh}}\nolimits}
\def\NL{\mathop{N\!L}\nolimits}
\def\CH{\mathop{\mathrm{CH}}\nolimits}
\def\proetale{{pro\text{-}\acute{e}tale}}
\def\etale{{\acute{e}tale}}
\def\QCoh{\mathit{QCoh}}
\def\Ker{\mathop{\mathrm{Ker}}}
\def\Im{\mathop{\mathrm{Im}}}
\def\Coker{\mathop{\mathrm{Coker}}}
\def\Coim{\mathop{\mathrm{Coim}}}

% Boxtimes
%
\DeclareMathSymbol{\boxtimes}{\mathbin}{AMSa}{"02}

%
% Macros for moduli stacks/spaces
%
\def\QCohstack{\mathcal{QC}\!\mathit{oh}}
\def\Cohstack{\mathcal{C}\!\mathit{oh}}
\def\Spacesstack{\mathcal{S}\!\mathit{paces}}
\def\Quotfunctor{\mathrm{Quot}}
\def\Hilbfunctor{\mathrm{Hilb}}
\def\Curvesstack{\mathcal{C}\!\mathit{urves}}
\def\Polarizedstack{\mathcal{P}\!\mathit{olarized}}
\def\Complexesstack{\mathcal{C}\!\mathit{omplexes}}
% \Pic is the operator that assigns to X its picard group, usage \Pic(X)
% \Picardstack_{X/B} denotes the Picard stack of X over B
% \Picardfunctor_{X/B} denotes the Picard functor of X over B
\def\Pic{\mathop{\mathrm{Pic}}\nolimits}
\def\Picardstack{\mathcal{P}\!\mathit{ic}}
\def\Picardfunctor{\mathrm{Pic}}
\def\Deformationcategory{\mathcal{D}\!\mathit{ef}}
\usepackage{cleveref}
\crefname{originalreference}{original reference}{original references}
\Crefname{originalreference}{Original reference}{Original references}
\makeatletter
\newlabel{perfect-lemma-relative-descend-homomorphisms}{{0DI7}{source}{}{sourcepointer-1}{}}
\newlabel{perfect-lemma-relative-descend-homomorphisms@cref}{{[originalreference][0][]0DI7}{[1][0][]source}{}{}{}}
\makeatother
\crefname{originalreference}{original reference}{original references}
\Crefname{originalreference}{Original reference}{Original references}
\makeatletter
\newlabel{situation-relative-descent}{{0DKU}{source}{}{sourcepointer-2}{}}
\newlabel{situation-relative-descent@cref}{{[originalreference][0][]0DKU}{[1][0][]source}{}{}{}}
\makeatother
\crefname{originalreference}{original reference}{original references}
\Crefname{originalreference}{Original reference}{Original references}
\makeatletter
\newlabel{spaces-perfect-lemma-derived-quasi-coherent-small-etale-site}{{071Q}{source}{}{sourcepointer-3}{}}
\newlabel{spaces-perfect-lemma-derived-quasi-coherent-small-etale-site@cref}{{[originalreference][0][]071Q}{[1][0][]source}{}{}{}}
\makeatother
\crefname{originalreference}{original reference}{original references}
\Crefname{originalreference}{Original reference}{Original references}
\makeatletter
\newlabel{spaces-perfect-lemma-descend-pseudo-coherent}{{08HE}{source}{}{sourcepointer-4}{}}
\newlabel{spaces-perfect-lemma-descend-pseudo-coherent@cref}{{[originalreference][0][]08HE}{[1][0][]source}{}{}{}}
\makeatother
\crefname{originalreference}{original reference}{original references}
\Crefname{originalreference}{Original reference}{Original references}
\makeatletter
\newlabel{spaces-perfect-lemma-induction-principle}{{08GP}{source}{}{sourcepointer-5}{}}
\newlabel{spaces-perfect-lemma-induction-principle@cref}{{[originalreference][0][]08GP}{[1][0][]source}{}{}{}}
\makeatother
\crefname{originalreference}{original reference}{original references}
\Crefname{originalreference}{Original reference}{Original references}
\makeatletter
\newlabel{spaces-perfect-lemma-mayer-vietoris-hom}{{08H9}{source}{}{sourcepointer-6}{}}
\newlabel{spaces-perfect-lemma-mayer-vietoris-hom@cref}{{[originalreference][0][]08H9}{[1][0][]source}{}{}{}}
\makeatother
\crefname{originalreference}{original reference}{original references}
\Crefname{originalreference}{Original reference}{Original references}
\makeatletter
\newlabel{spaces-perfect-lemma-pseudo-coherent}{{08JL}{source}{}{sourcepointer-7}{}}
\newlabel{spaces-perfect-lemma-pseudo-coherent@cref}{{[originalreference][0][]08JL}{[1][0][]source}{}{}{}}
\makeatother
\crefname{originalreference}{original reference}{original references}
\Crefname{originalreference}{Original reference}{Original references}
\makeatletter
\newlabel{spaces-perfect-lemma-tor-dimension-rel}{{0DK7}{source}{}{sourcepointer-8}{}}
\newlabel{spaces-perfect-lemma-tor-dimension-rel@cref}{{[originalreference][0][]0DK7}{[1][0][]source}{}{}{}}
\makeatother
\title{Descending homomorphisms between derived objects along a relative limit}
\author{The Stacks Project Authors}
\date{stacks-project:spaces-more-morphisms.tex:local-snapshot-20261011; GNU FDL 1.2 or later}
\begin{document}
\maketitle
\noindent\textit{Editorial note.} Counted span is the author's own proof environment, source lines 12451-12496, immediately following the statement at 12434-12449.\par\medskip
\noindent\textit{Editorial note.} Seven cross-chapter pointers (Derived Categories of Spaces/Schemes) plus the chapter's own situation are recorded in assembly.external\_references.\par\medskip
\noindent\textit{Editorial note.} The colimit comparison map is displayed verbatim in the statement span.\par\medskip
\noindent\textit{Allowed standard prerequisites (editorial).}\begin{itemize}
\item The situation of relative descent. chapter situation situation-relative-descent, cited in the statement. Fixes the limit presentation X = lim X\_i over X\_0 and the pullbacks K\_i, L\_i whose Hom sets are compared.
\item The schemes-level descent statement. Derived Categories of Schemes, lemma perfect-lemma-relative-descend-homomorphisms, cited in the proof. Supplies the argument in the scheme case, which is transferred to algebraic spaces here.
\item Induction principle and Mayer-Vietoris for derived categories of spaces. Derived Categories of Spaces, lemmas spaces-perfect-lemma-induction-principle and spaces-perfect-lemma-mayer-vietoris-hom, cited in the proof. Reduce the verification to affine pieces where the comparison can be made.
\item Descending pseudo-coherence and tor dimension. Derived Categories of Spaces, lemmas spaces-perfect-lemma-descend-pseudo-coherent and spaces-perfect-lemma-tor-dimension-rel, cited in the proof. Keep the relevant finiteness conditions along the descent.
\item The small etale site comparison. Derived Categories of Spaces, lemma spaces-perfect-lemma-derived-quasi-coherent-small-etale-site, cited in the proof. Moves the computation between the derived category of the space and that of its small etale site.
\end{itemize}\medskip
\ifdefined\hypertarget\hypertarget{sourcepointer-1}{}\fi
\noindent\textit{Source reference.} Original reference 0DI7: focused-seeds/geometry/S9-spaces-relative-descend-homomorphisms-spaces/ref-tags.txt, source lines 5--5. Navigation pointer only: the counted proof cites this statement by \textbackslash{}ref as a standing input; the deduction the counted proof has to make is carried out inside the counted span, and no part of that deduction is delegated to this pointer.\par\medskip
\ifdefined\hypertarget\hypertarget{sourcepointer-2}{}\fi
\noindent\textit{Source reference.} Original reference 0DKU: sources/stacks-spaces-more-morphisms/pinned-original-source.tex, source lines 12421--12421. Navigation pointer only: the counted proof cites this statement by \textbackslash{}ref as a standing input; the deduction the counted proof has to make is carried out inside the counted span, and no part of that deduction is delegated to this pointer.\par\medskip
\ifdefined\hypertarget\hypertarget{sourcepointer-3}{}\fi
\noindent\textit{Source reference.} Original reference 071Q: focused-seeds/geometry/S9-spaces-relative-descend-homomorphisms-spaces/ref-tags.txt, source lines 6--6. Navigation pointer only: the counted proof cites this statement by \textbackslash{}ref as a standing input; the deduction the counted proof has to make is carried out inside the counted span, and no part of that deduction is delegated to this pointer.\par\medskip
\ifdefined\hypertarget\hypertarget{sourcepointer-4}{}\fi
\noindent\textit{Source reference.} Original reference 08HE: focused-seeds/geometry/S9-spaces-relative-descend-homomorphisms-spaces/ref-tags.txt, source lines 7--7. Navigation pointer only: the counted proof cites this statement by \textbackslash{}ref as a standing input; the deduction the counted proof has to make is carried out inside the counted span, and no part of that deduction is delegated to this pointer.\par\medskip
\ifdefined\hypertarget\hypertarget{sourcepointer-5}{}\fi
\noindent\textit{Source reference.} Original reference 08GP: focused-seeds/geometry/S9-spaces-relative-descend-homomorphisms-spaces/ref-tags.txt, source lines 8--8. Navigation pointer only: the counted proof cites this statement by \textbackslash{}ref as a standing input; the deduction the counted proof has to make is carried out inside the counted span, and no part of that deduction is delegated to this pointer.\par\medskip
\ifdefined\hypertarget\hypertarget{sourcepointer-6}{}\fi
\noindent\textit{Source reference.} Original reference 08H9: focused-seeds/geometry/S9-spaces-relative-descend-homomorphisms-spaces/ref-tags.txt, source lines 9--9. Navigation pointer only: the counted proof cites this statement by \textbackslash{}ref as a standing input; the deduction the counted proof has to make is carried out inside the counted span, and no part of that deduction is delegated to this pointer.\par\medskip
\ifdefined\hypertarget\hypertarget{sourcepointer-7}{}\fi
\noindent\textit{Source reference.} Original reference 08JL: focused-seeds/geometry/S9-spaces-relative-descend-homomorphisms-spaces/ref-tags.txt, source lines 10--10. Navigation pointer only: the counted proof cites this statement by \textbackslash{}ref as a standing input; the deduction the counted proof has to make is carried out inside the counted span, and no part of that deduction is delegated to this pointer.\par\medskip
\ifdefined\hypertarget\hypertarget{sourcepointer-8}{}\fi
\noindent\textit{Source reference.} Original reference 0DK7: focused-seeds/geometry/S9-spaces-relative-descend-homomorphisms-spaces/ref-tags.txt, source lines 11--11. Navigation pointer only: the counted proof cites this statement by \textbackslash{}ref as a standing input; the deduction the counted proof has to make is carried out inside the counted span, and no part of that deduction is delegated to this pointer.\par\medskip
\begin{lemma}
\label{lemma-relative-descend-homomorphisms}
In Situation \ref{situation-relative-descent}.
Let $K_0$ and $L_0$ be objects of $D(\mathcal{O}_{X_0})$.
Set $K_i = Lf_{i0}^*K_0$ and $L_i = Lf_{i0}^*L_0$ for $i \geq 0$
and set $K = Lf_0^*K_0$ and $L = Lf_0^*L_0$. Then the map
$$
\colim_{i \geq 0} \Hom_{D(\mathcal{O}_{X_i})}(K_i, L_i)
\longrightarrow
\Hom_{D(\mathcal{O}_X)}(K, L)
$$
is an isomorphism if $K_0$ is pseudo-coherent and
$L_0 \in D_\QCoh(\mathcal{O}_{X_0})$ has (locally)
finite tor dimension as an object of
$D((X_0 \to Y_0)^{-1}\mathcal{O}_{Y_0})$
\end{lemma}

\begin{proof}
For every quasi-compact and quasi-separated object $U_0$ of
$(X_0)_{spaces, \etale}$ consider the condition $P$ that
$$
\colim_{i \geq 0} \Hom_{D(\mathcal{O}_{U_i})}(K_i|_{U_i}, L_i|_{U_i})
\longrightarrow
\Hom_{D(\mathcal{O}_U)}(K|_U, L|_U)
$$
is an isomorphism where $U = X \times_{X_0} U_0$ and
$U_i = X_i \times_{X_0} U_0$. We will prove $P$ holds for each $U_0$.

\medskip\noindent
Suppose that $(U_0 \subset W_0, V_0 \to W_0)$ is an elementary
distinguished square in $(X_0)_{spaces, \etale}$
and $P$ holds for $U_0, V_0, U_0 \times_{W_0} V_0$.
Then $P$ holds for $W_0$ by Mayer-Vietoris
for hom in the derived category, see Derived Categories of Spaces,
Lemma \ref{spaces-perfect-lemma-mayer-vietoris-hom}.

\medskip\noindent
We first consider $U_0 = W_0 \times_{Y_0} X_0$ with $W_0$ a
quasi-compact and quasi-separated object of $(Y_0)_{spaces, \etale}$.
By the induction principle of Derived Categories of Spaces,
Lemma \ref{spaces-perfect-lemma-induction-principle}
applied to these $W_0$ and the previous paragraph,
we find that it is enough to prove
$P$ for $U_0 = W_0 \times_{Y_0} X_0$ with $W_0$ affine.
In other words, we have reduced to the case where $Y_0$ is affine.
Next, we apply the induction principle again, this time to all
quasi-compact and quasi-separated opens of $X_0$, to reduce to the
case where $X_0$ is affine as well.

\medskip\noindent
If $X_0$ and $Y_0$ are affine, then we are back in the case
of schemes which is proved in
Derived Categories of Schemes, Lemma
\ref{perfect-lemma-relative-descend-homomorphisms}.
The reader may use
Derived Categories of Spaces, Lemmas
\ref{spaces-perfect-lemma-pseudo-coherent},
\ref{spaces-perfect-lemma-derived-quasi-coherent-small-etale-site},
\ref{spaces-perfect-lemma-descend-pseudo-coherent}, and
\ref{spaces-perfect-lemma-tor-dimension-rel}
to accomplish the translation of the statement into a statement
involving only schemes and derived categories of modules on schemes.
\end{proof}

\end{document}
